\section{Auditoría de fuentes y eje del curso: por qué el audio es una prueba de esfuerzo para el Transformer}

Esta clase no se limita a trasladar el Transformer de NLP al sonido; plantea una pregunta de diseño más básica: el sonido es continuo, tiene una frecuencia de muestreo alta y abarca varias escalas de tiempo; ¿cómo debe representarse primero y qué objetivo de predicción debe recibir después? Prateek Verma enlaza tres trabajos que abarcan raw waveform synthesis, self-supervised representation learning y large-scale audio understanding. Las tres líneas parecen corresponder por separado a generación, representación y clasificación, pero en realidad comparten la misma cadena de ingeniería: la \term{representation} determina la longitud de la secuencia y la pérdida de información, el \term{objective} determina qué conserva el modelo y la \term{metric} determina qué es lo que realmente demostramos.

La página oficial del curso registra esta clase el 2021-11-29, y Stanford Online no subió la grabación hasta el 2022-07-18. El artículo Audio Transformers que corresponde al video estaba en su versión v1 de 2021 al momento de la clase; arXiv publicó después, en 2025, una versión v2. Estas notas usan solo la v1 para explicar los experimentos de la clase; no se puede proyectar hacia atrás una revisión de cuatro años después como si fuera evidencia con la que el expositor ya contaba. El video no publicó un PDF independiente de las diapositivas, por lo que, de los 1,449 puntos muestreados de la grabación oficial en 1080p, se seleccionaron manualmente 47 teaching states, y se usaron los 981 subtítulos oficiales hechos por personas para recuperar las explicaciones orales del docente.

\lecturefigure{slide-01-title-understanding-to-synthesis.jpg}{El título de la clase resume el eje así: de language modeling a understanding y, después, a synthesis.}{Video oficial de Stanford Online, 00:00:24--00:00:47.}

Los tres verbos de la diapositiva de título no son una clasificación publicitaria. \emph{Synthesis} exige generar formas de onda de alta frecuencia; \emph{understanding} exige comprimir una señal larga en una representación útil para las etiquetas; \emph{language modeling} exige convertir primero el sonido continuo en símbolos discretos predecibles. Al leer lo que sigue conviene preguntarse continuamente: ¿el sistema actual procesa un sample, un spectrum, un embedding o un code index? Si esa capa no queda clara, el llamado «Audio Transformer» es solo un nombre ambiguo.

\teachervoice{Desde el inicio, Verma llama a esta clase ``buffet style'': no va a demostrar que una arquitectura unificada gana siempre, sino que usará tres artículos para mostrar varias ideas combinables. Por eso estas notas se organizan por problemas comunes y no dividen los capítulos mecánicamente según los resúmenes de los artículos.}

\subsection{Cómo forman los tres artículos una cadena de diseño}

El primer trabajo hace next-sample prediction sobre raw samples de 8 bits y compara WaveNet con un causal Transformer; el segundo convierte primero el audio continuo en codebook tokens y luego compara contrastive learning con next-code prediction; el tercero hace que el Transformer aprenda directamente representaciones para clasificación a partir de la raw waveform, y añade pooling o una Haar wavelet sobre el embedding intermedio. Los puntos de conexión entre los tres trabajos son los siguientes:

\lecturefigure{slide-02-talk-roadmap.jpg}{Hoja de ruta de la clase: fundamentos de representación, raw synthesis, representation learning con tokens discretos y raw-waveform understanding.}{Video oficial de Stanford Online, 00:01:08--00:01:45.}

\begin{knowledgebox}{Lectura de la figura: no confundir las tres rutas con «el mismo Transformer»}
La salida de raw synthesis es el siguiente sample cuantizado; la salida de representation learning puede ser una code prediction o un linear probe posterior; la salida de audio understanding son sound-event scores multietiqueta. Comparten attention, pero difieren en la granularidad de la entrada, la señal de supervisión, la longitud de la secuencia y la métrica de evaluación. La unidad de comparación correcta no es el nombre del modelo, sino la terna «representación--objetivo--métrica».
\end{knowledgebox}

\begin{warningbox}{Límite temporal de la clase}
En 2021 el expositor describió el Transformer como «la tendencia del momento» en varios campos y usó «Adieu WaveNet?» y «Adieu Convolutions» como títulos provocadores. Estos títulos solo explican los tres grupos de experimentos controlados de entonces; de ellos no se puede concluir que en 2026 todos los sistemas de voz, música, codec o generación deban abandonar las convoluciones, los modelos de espacio de estados o los modelos de difusión.
\end{warningbox}

\subsection{Con qué grado de solidez debe leerse la evidencia}

Lo que sigue incluye hechos algebraicos, pero también benchmark tables, visualizaciones de filtros y juicios del expositor sobre el futuro. Para no presentarlos como conclusiones de la misma solidez, esta clase adopta cuatro niveles de evidencia: si algo es realizable en cuanto a estructura, si mejora en experimentos controlados, si la representación interna es interpretable y si la conclusión puede extrapolarse a la calidad perceptual o a otras tareas. Ningún nivel sustituye automáticamente al siguiente.

\begin{table}[H]
\centering
\small
\caption{Evidencia frecuente en esta clase y conclusiones que puede sustentar.}
\begin{tabularx}{\textwidth}{L{0.21\textwidth}L{0.29\textwidth}X}
\toprule
Evidencia & Ejemplo & Conclusión razonable \\
\midrule
Estructura y complejidad & $n=f_sT$, attention es $O(n^2)$ & Explica por qué el raw audio requiere ventanas locales, compresión o estructuras dispersas. \\
Métricas controladas & top-5 accuracy o mAP con los mismos datos & Describe el desempeño de predicción/clasificación en esa configuración; no equivale a calidad de audio subjetiva ni a una ventaja arquitectónica universal. \\
Sondas de representación & Linear probe, análisis de respuesta en frecuencia de los filtros aprendidos & Indica que la representación contiene información utilizable o que aparecen estructuras semejantes al procesamiento de señales clásico; no demuestra que una sola neuron tenga un significado único. \\
Extrapolación y juicio & «Más datos reducirán el gap», «el Transformer no es el punto final» & Son hipótesis o líneas de investigación; deben separarse de los resultados medidos. \\
\bottomrule
\end{tabularx}
\end{table}

\subsection{Resumen de la sección}

El objeto de esta clase no es un modelo único, sino el espacio de diseño del modelado de audio. La fecha de la clase, la fecha de publicación del video y la fecha de revisión del artículo deben mantenerse separadas; los resultados de los tres trabajos también deben interpretarse dentro de sus respectivas representaciones de entrada, funciones objetivo y alcances de las métricas.
